\scanpage{682}

% 682-01, upper heading
$8^\circ)$ Retour sur le cas universel, et l'hom
\[
 M^{\unsure{\natural|\sharp}}\longrightarrow
 M_{\rho,\sigma}^{\unsure{\natural|\sharp}}/\mu_\rho
\]

% 682-03, upper and middle text
Nous supposons maintenant que $\unsure{\mathcal G}$ correspond au cas universel
$\mathrm{GL}(2,\mathbb Z)^{\wedge}$, avec un
$\unsure{\mathcal G}^{\unsure{\natural|\sharp}}$ quelconque. On va noter le cas
universel par un indice $0$ :
$\unsure{\mathcal G}_0,\ \unsure{\mathcal G}_0^+,\
\unsure{\mathcal G}_0^{\unsure{\natural|\sharp}}$ etc., et
$\unsure{\mathcal G}',\ \unsure{\mathcal G}'^{\unsure{\natural|\sharp}}$
etc. \unsure{les|le} \unsure{ci-devant} $\unsure{\mathcal G},\
\unsure{\mathcal G}^{\unsure{\natural|\sharp}}$ etc.

% 682-04, left margin at middle
\margin{$\unsure{\swarrow}\ Z_{\rho_0,\sigma_0}$ etc $\ldots$}
\margin{On écrira $\unsure{\mathcal G}_0^+$ au lieu de
$\unsure{\mathcal G}_0^{+\wedge}$ $\ldots$}

% 682-06, middle-lower text and display
Je pose
\[
\begin{array}{rl}
(107)\qquad M^{\unsure{\natural|\sharp}}&\subset M\\[-2pt]
\Vert&\\[-6pt]
\displaystyle
\left\{u\in M\ \middle|\
\begin{array}{ll}
\text{a) }u(\rho)\text{ est }\unsure{\mathcal G}^{\unsure{\natural|\sharp}}\text{-conjugué à}
 &\varepsilon_3[\tau](\rho)\\
\text{b) }u(\sigma)\quad\rule{4.2cm}{0.3pt}&\varepsilon_2[\tau](\sigma)\\
\text{c) }u(\varepsilon_0)\quad\rule{4.2cm}{0.3pt}&\varepsilon_0^\mu\\
\text{d) }u\text{ stabilise }\unsure{M}^{\unsure{\natural|\sharp}}&
\end{array}\right\}&
\end{array}
\]
\margin{\uline{NB} $M^{\unsure{\natural|\sharp}}\unsure{\supset|\ni}\mu$}

% 682-09, lower text and display
où bien sûr
\[
\mu=\chi(u),\quad
\varepsilon_2=\varepsilon_2(u)=\varepsilon_2(\mu),\quad
\varepsilon_3=\varepsilon_3(u)=\varepsilon_3(\mu).
\]
Je dis que $M^{\unsure{\natural|\sharp}}$ est un \uline{sous-groupe} de $M$.
On peut le voir en \unsure{reprenant} les calculs de la p.\ $629$ ff.
Je \unsure{préfère} procéder ainsi.

\margin{681}

% Printed archive footer.
\begin{center}
681\\
\texttt{https://grothendieck.umontpellier.fr}\\
Université de Montpellier, UMR CNRS 5149, Case courrier 051, Place Eugène Bataillon\\
34095 Montpellier cedex 5 - France
\end{center}


\scanpage{683}
\mspage{644}

% 683-02, upper text and display
On a vu qu'on a
\[
(108)\qquad
\underset{\deleted{\unsure{\sharp}}}{Z_{\rho_0,\sigma_0}}
\hookrightarrow M[0]\inserted{{}=M(0)\cdot L_0^{\unsure{\mathcal G}}}
\simeq \widetilde M_{0,3}[0]
=\widetilde M_{0,3}(0)\cdot L_0^{\unsure{\mathcal G}_{0,3}}
\]
inclusion qui est une égalité si
$\unsure{\mathcal G}_0^{\unsure{\natural|\sharp}}=\unsure{\mathcal G}_0^+$,
et qui identifie $Z_{\rho,\sigma}$ à
\[
M^{!}[0]\simeq M(0)\cdot L_0^\pi
\quad\text{si}\quad
\unsure{\mathcal G}_0^{\unsure{\natural|\sharp}}=\pi\unsure{:=}\Pi_0.
\]

% 683-08, bracketed passage
[ --- en fait, si
$\unsure{\mathcal G}_0^{\unsure{\natural|\sharp}}\subset\Pi_0$,
alors on trouve
\[
\begin{split}
(109)\quad Z_{\rho,\sigma}=\left\{
u=\underset{\substack{\rotatebox{-90}{$\in$}\\M(0)}}{u_0}\cdot
\underset{\substack{\rotatebox{-90}{$\in$}\\\deleted{\unsure{\mathcal G}_0}L_0^\pi}}{\ell}
\ \middle|\
\begin{array}{l}
\text{a) }u(\rho)=\alpha(\rho)\quad\text{avec }\alpha\in
\unsure{\mathcal G}_{\unsure{\bar\tau}}^{\unsure{\natural|\sharp}}\\
\text{b) }u(\sigma)=\beta(\sigma)\quad\text{avec }\beta\in
\unsure{\mathcal G}_{\unsure{\bar\tau}}^{\unsure{\natural|\sharp}}\\
\text{d) }u\unsure{\varepsilon|\in}\text{ stabilise }
\unsure{\mathcal G}_0^{\unsure{\natural|\sharp}}
\end{array}\right\}
\end{split}
\]
\margin{cette explicitation ne \unsure{semble} servir de rien.}

% 683-12, middle-lower text with cancellations and additions
et on voit que $\alpha,\beta$ dans a) b) sont déterminés par
$u=u_0\ell$ de façon unique
(\deleted{dès que} $\unsure{\overline{\mathrm m}}$
\unsure{si on|l'on} suppose seulement $u\in M^{!}[0]$ sans plus,
et \inserted{\unsure{qu'on exige}}
\deleted{\unsure{$\alpha\in$|avec} $\unsure{\mathcal G}_0^{\unsure{\natural}}$}
$\alpha,\unsure{\beta}\in\unsure{\Pi_{0\bar\tau}}$,
\deleted{\unsure{peut|Bref}}
On sait que c'est un \unsure{s groupe|ss groupe}
\inserted{de $\operatorname{Aut}(\unsure{\mathcal G}_0^+)\times
\widehat{\mathbb Z}^{*}\times\mu\times\mu$
--- et en fait, de $\operatorname{Aut}(\unsure{\mathcal G}_0^+)$}
(prop.\ p.\ $629$).
\deleted{\unsure{Considérons}}

% 683-19, lower paragraph
Comme ce groupe opère sur \deleted{$M$ et}
$\unsure{\mathcal G}$ via $\widetilde M_{0,3}$,
\inserted{$=\text{ via }\widetilde M_{0,3}^{!}[0]$}
il s'envoie dans $\widetilde M_{0,3}^{!}[0]$ par un hom de groupes,
donc dans \deleted{$\unsure{\widetilde M_{0,3}[0]}$}
$M[0]^{!}\simeq\widetilde M_{0,3}[0]$ par un hom de groupes. ]

\margin{682}

% Printed archive footer.
\begin{center}
682\\
\texttt{https://grothendieck.umontpellier.fr}\\
Université de Montpellier, UMR CNRS 5149, Case courrier 051, Place Eugène Bataillon\\
34095 Montpellier cedex 5 - France
\end{center}


\scanpage{684}

% 684-01, upper text and tuple
\deleted{\unsure{L'hom}} Considérons le système
\[
\begin{split}
(110)\quad\bigl(&\inserted{\unsure{\mathcal G}_0^{+\wedge},}
Z_{\rho_0,\sigma_0},\quad
L_0^{\unsure{\natural|\sharp}}\subset Z_{\rho_0,\sigma_0},\quad
Z_{\rho_0,\sigma_0}\longrightarrow
\operatorname{Aut}(\unsure{\mathcal G}_0^{\unsure{+\wedge}}),\quad
L_0^{\unsure{\natural|\sharp}}\hookrightarrow\unsure{\mathcal G}_0^{+\wedge}\bigr)
\end{split}
\]
(cf ``digression'' plus bas $9^\circ$) qui donne naissance à
\[
(110')\qquad
G_{\rho_0,\sigma_0},\quad
\unsure{\mathcal G}_0^{+\wedge}\subset G_{\rho_0,\sigma_0},\quad
Z_{\rho_0,\sigma_0}\hookrightarrow G_{\rho_0,\sigma_0},
\]
et considérons le système analogue \inserted{à $(110^{*})$}
\[
\begin{split}
(111^{\unsure{*|\sharp}})\quad\bigl(&\inserted{\unsure{\mathcal G}_0^+,}
\underset{\substack{\Vert\\M(0)\cdot L_0^\pi}}{M[0]^{\unsure{!|\natural}}},\quad
\underset{\text{(\unsure{engendré} par }\varepsilon_0\text{)}}{L_0^{\unsure{\mathcal G}}}
\subset M[0]^{\unsure{!|\natural}},\quad
\unsure{\deleted{M[0]}}^{\unsure{!|\natural}}\longrightarrow
\operatorname{Aut}(\unsure{\mathcal G}_0^+),\quad
L_0^{\unsure{\mathcal G}}\hookrightarrow\unsure{\mathcal G}_0^+\bigr)
\end{split}
\]
qui donne naissance (loc.\ cit ($9^\circ$)) à la situation
\[
(111')\qquad
\bigl(M,\quad\unsure{\mathcal G}_0^+\subset M,\quad
M[0]'\hookrightarrow\unsure{\mathcal G}_0^+\bigr).
\]

% 684-14, lower paragraph
Or grâce à $(108)$, on a un hom de $(\unsure{110|109|119})$ dans $(111)$,
correspondant à l'identité sur $\unsure{\mathcal G}_0^+$,
$(108)$ sur $Z_{\rho_0,\sigma_0}$, l'inclusion can.\ :
\deleted{\unsure{sur}}
$L_0^{\unsure{\natural|\sharp}}\subset L_0^{\unsure{\mathcal G}}$.
Il en résulte un hom de $(110')$ dans $(111')$, en particulier un hom
canonique de suites exactes

\margin{683}

% Printed archive footer.
\begin{center}
683\\
\texttt{https://grothendieck.umontpellier.fr}\\
Université de Montpellier, UMR CNRS 5149, Case courrier 051, Place Eugène Bataillon\\
34095 Montpellier cedex 5 - France
\end{center}


\scanpage{685}
\mspage{645}

% 685-02, upper commutative diagram
% (112): the first \mathcal G_0^+ seems written over an earlier G_\rho (unresolved).
\begin{samepage}
(112)\quad
\diagram{\begin{tikzcd}[column sep=2.1em,row sep=2.6em]
1 \arrow[r,no head] &
\unsure{\mathcal G}_0^+ \arrow[r] \arrow[d,"\unsure{1}"] &
G_{\rho_0,\sigma_0} \arrow[r] \arrow[d,hook] &
\gamma_{\rho_0,\sigma_0} \arrow[r] \arrow[d,hook,dotted] \arrow[d] & 1 \\
1 \arrow[r,no head] &
\unsure{\mathcal G}_0^+ \arrow[r,no head] & M \arrow[r] &
\widetilde\Gamma_{\mathbb Q} \arrow[r,no head] & 1
\end{tikzcd}}
\end{samepage}

% 685-03, explanatory diagram and text
où
\diagram{\begin{tikzcd}[column sep=4em,row sep=1em]
\gamma_{\rho_0,\sigma_0} \arrow[r] \arrow[d,phantom,"\unsure{\wr}" description] &
\widetilde\Gamma_{\mathbb Q} \arrow[d,phantom,"\unsure{\wr}" description] \\
Z_{\rho_0,\sigma_0}/L_0^{\unsure{\mu}} &
M[0]/L_0^{\unsure{\mathcal G}}
\end{tikzcd}}
est \inserted{\unsure{celui}} induit par
$(\unsure{108})$, et s'identifie à :
\[
(113)\qquad Z_{\rho_0,\sigma_0}/L_0^{\unsure{\natural|\sharp}}
\longrightarrow M[0]/L_0^{\unsure{\mathcal G}}
\]

% 685-07, middle text
Par définition $\overline{\mathrm m}$ de
$L_0^{\unsure{\natural|\sharp}}$ comme étant
$Z_{\rho_0,\sigma_0}\cap\unsure{\mathcal G}_0^+$
\deleted{\unsure{et|ce}}, on voit que $(113)$ est inj.\ donc
$G_{\rho_0,\sigma_0}\longrightarrow M$ est tjrs injectif
(c'est un iso.\ si
$\unsure{\mathcal G}_0^{\unsure{\natural|\sharp}}\supset\Pi_0$
$\ldots$).
On déduit par comparaison de $(112)$ un hom de suites exactes

% 685-09, lower commutative diagram
\begin{samepage}
(114)\quad
\diagram{\begin{tikzcd}[column sep=2.1em,row sep=2.6em]
1 \arrow[r,no head] &
\unsure{\mathcal G}_0^{\unsure{\natural|\sharp}} \arrow[r] \arrow[d,hook] &
G_{\rho_0,\sigma_0}^{\unsure{\natural|\sharp}} \arrow[r] \arrow[d,hook] &
\gamma_{\rho_0,\sigma_0} \arrow[r,no head] \arrow[d,hook] & 1 \\
1 \arrow[r,no head] &
\unsure{\mathcal G}_0^{\unsure{+}} \arrow[r] &
M \arrow[r] &
\deleted{\unsure{\gamma}}\widetilde\Gamma_{\mathbb Q} \arrow[r,no head] & 1.
\end{tikzcd}}
\end{samepage}

% 685-10, bottom text
On voit alors immédiatement que l'image de
$G_{\rho_0,\sigma_0}^{\unsure{\natural|\sharp}}$ dans $M$ n'est autre que
$M^{\unsure{\natural|\sharp}}$ $(107)$. On a donc

\margin{684}

% Printed archive footer.
\begin{center}
684\\
\texttt{https://grothendieck.umontpellier.fr}\\
Université de Montpellier, UMR CNRS 5149, Case courrier 051, Place Eugène Bataillon\\
34095 Montpellier cedex 5 - France
\end{center}


\scanpage{686}

\[
(115)\qquad
\left\{
\begin{array}{ll}
\underset{\substack{\text{\unsure{défini dans}}\\(107)}}{M^{\unsure{\natural|\sharp}}}
  \simeq G_{\rho_0,\sigma_0}^{\unsure{\natural|\sharp}},\\[5pt]
\gamma_{\rho_0,\sigma_0}\hookrightarrow\widetilde{\Gamma}_{\mathbb Q}
  \deleted{\unsure{\Leftarrow}\{u\in\widetilde{\Gamma}_{\mathbb Q}}
\end{array}
\right.
\]

\margin{
\diagram{
\begin{tikzcd}[row sep=small]
\text{\uline{à vérifier}}\arrow[d]\\
\text{\uline{NB}\quad On a}
\end{tikzcd}
}
\[
(116)\quad
\left\{
\begin{array}{l}
\gamma_{\rho_0,\sigma_0}
 \deleted{\unsure{=\Gamma_0}}
 \inserted{\unsure{=\Gamma_0^{\natural}|\Gamma_0^{\sharp}}}
 =\left\{u\in\widetilde{\Gamma}_{\mathbb Q}\ \middle|\
 \begin{array}{l}
\text{le relèvement }u_0\\
\text{de }u\text{ en }u_0\in M(0)\\
\text{satisfait les}\\
\text{conditions a) b) d)}\\
\text{de (107)}
\end{array}\right\}\\
\text{(\unsure{notations}}\\
\gamma_0=\Gamma_0^{\unsure{\natural|\sharp}},
\ \text{où }\Gamma_0\text{ dénote }\widetilde{\Gamma}_{\mathbb Q}\text{)}
\end{array}
\right.
\]
}

qui \unsure{montre} en $\overline{\mathrm m}$ temps dans la
\unsure{foulée} que $M^{\unsure{\natural|\sharp}}$ est bien un
sous-groupe de $M$.

On a maintenant
\[
(\unsure{117})\qquad
M_{\rho_0,\sigma_0}^{\unsure{\natural|\sharp}}
\overset{\unsure{\mathrm{dfn}|\mathrm{déf}}}{\simeq}
\unsure{N|M}_{\rho_0}^{\unsure{\natural|\sharp}}
\times_{\gamma_{\rho_0,\sigma_0}}
N_{\sigma_0}^{\unsure{\natural|\sharp}}
\times_{\unsure{\gamma_{\rho_0,\sigma_0}}}
\underbrace{G_{\rho_0,\sigma_0}^{\unsure{\natural|\sharp}}}
 _{\simeq M^{\unsure{\natural|\sharp}}},
\]
où $G_{\rho_0,\sigma_0}^{\unsure{\natural|\sharp}}\simeq
M^{\unsure{\natural|\sharp}}$ \unsure{donc|donne}
\[
\begin{aligned}
N_{\rho_0}^{\unsure{\natural|\sharp}}
 &\unsure{\simeq|=}\left\{u\in M^{\unsure{\natural|\sharp}}\mid
 u(\rho)=\rho^{\varepsilon_3(u)}\right\}\\
N_{\sigma_0}^{\unsure{\natural|\sharp}}
 &\unsure{\simeq|=}\left\{u\in M^{\unsure{\natural|\sharp}}\mid
 u(\sigma)\unsure{\equiv|=}\sigma^{\varepsilon_2(u)}\right\}
\end{aligned}
\]
et finalement on a
\[
(118)\qquad
\left\{
\begin{array}{l}
M_{\rho_0,\sigma_0}^{\unsure{\natural|\sharp}}
\simeq
N_{M^{\unsure{\natural|\sharp}}}(\rho)
\times_{\deleted{\unsure{\gamma_{\rho_0,\sigma_0}}}
 \inserted{\Gamma_0^{\unsure{\natural|\sharp}}}}
N_{M^{\unsure{\natural|\sharp}}}(\sigma)
\times_{\deleted{\unsure{\gamma_{\rho_0,\sigma_0}}}
 \inserted{\Gamma_0^{\unsure{\natural|\sharp}}}}
M^{\unsure{\natural|\sharp}}.\\[5pt]
\text{où }\Gamma_0^{\unsure{\natural|\sharp}}\unsure{\ne|=}\gamma_0
\simeq\gamma_{\rho_0,\sigma_0}\subset\widetilde{\Gamma}_{\mathbb Q}
\quad\text{\unsure{explicité} dans (116)},
\end{array}
\right.
\]
\uuline{d'où} une \deleted{isom. de} \unsure{suites|suite}
\unsure{exactes|exacte}
\begin{samepage}
\[
(119)
\]
\diagram{
\begin{tikzcd}[column sep=small]
1\arrow[r]&
\overbrace{SZ_{\rho_0}\times SZ_{\sigma_0}}^{\inserted{L_\rho\times L_\sigma}}
\arrow[r]&
M_{\rho_0,\sigma_0}^{\unsure{\natural|\sharp}}
\arrow[r]\arrow[d,"\inserted{\unsure{\natural|\sharp}}"']&
\deleted{\unsure{\gamma}}M^{\unsure{\natural|\sharp}}\arrow[r]&1\\
&&{}&&
\end{tikzcd}
}
\end{samepage}
où le noyau $SZ_{\rho_0}\times SZ_{\sigma_0}$ est un groupe
\unsure{tout petit}, savoir

\margin{685}

% Printed archive footer.
\begin{center}
685\\
\texttt{https://grothendieck.umontpellier.fr}\\
Université de Montpellier, UMR CNRS 5149, Case courrier 051, Place Eugène Bataillon\\
34095 Montpellier cedex 5 - France
\end{center}


\scanpage{687}
\mspage{646}

\[
(120)\quad
\operatorname{Ker}\left(
M_{\rho_0,\sigma_0}^{\unsure{\natural|\sharp}}
\overset{\text{épi}}{\longrightarrow}M^{\unsure{\natural|\sharp}}
\right)
\simeq SZ_{\rho_0}\times SZ_{\sigma_0}
\simeq L_\rho\times L_\sigma
\simeq \mathbb Z/6\mathbb Z\times\mathbb Z/4\mathbb Z .
\]

Nous allons améliorer \unsure{encor|encore} la \deleted{\illegible{2}}\inserted{situation},
en découpant dans $M_{\rho_0,\sigma_0}^{\unsure{\natural|\sharp}}$
un groupe \unsure{encor|encore} plus petit, qui soit une ext. de
$M^{\unsure{\natural|\sharp}}$ par $\mu$.
\deleted{Pour ceci, on définit}
\[
\deleted{
N_{\unsure{\sigma_0|\rho_0}}^{\unsure{\natural|\sharp}!}
=N_{\sigma}^{\unsure{\natural|\sharp}}\cap M_0^!
\unsure{=|\cdot}M^{\unsure{\natural|\sharp}}\unsure{\cap}
}
\]

\unsure{(}

Notons que
\[
(121)\qquad
\begin{array}{c}
M_{\rho_0,\sigma_0}^{\unsure{\natural|\sharp}}
=N_{\rho_0}^{\unsure{\natural|\sharp}}
\times_{\Gamma_0^{\unsure{\natural|\sharp}}}
N_{\sigma_0}^{\unsure{\natural|\sharp}}
\times_{\Gamma_0^{\unsure{\natural|\sharp}}}
\deleted{\unsure{G_{\rho_0,\sigma_0}^{\natural}|G_{\rho_0,\sigma_0}^{\sharp}}}
\underbrace{\inserted{\unsure{G|M}_{\rho_0,\sigma_0}^{\unsure{\natural|\sharp}}}}
 _{M^{\unsure{\natural|\sharp}}}\\[5pt]
\cap\\[3pt]
N_{\rho_0}\times_{\Gamma_0}N_{\sigma_0}\times_{\Gamma_0}M .
\end{array}
\]
\margin{où $\Gamma_0\overset{\unsure{\mathrm{dfn}|\mathrm{déf}}}{=}
\widetilde{\Gamma}_{\mathbb Q}$, \unsure{et}}

\inserted{et}\ \unsure{où} ce dernier produit fibré n'est autre que la
quantité $M_{\rho_0,\sigma_0}^{(\unsure{\natural|\sharp}\,')}$
qui correspondrait \unsure{au|en} cas
\deleted{$\unsure{M_{\rho_0,\sigma_0}}$}
de $\unsure{\mathcal G}_0^{\unsure{\natural|\sharp}\,'}
=\unsure{\mathcal G}_0^+$ \unsure{tt|toute} entière
(noté $M_{\rho,\sigma}$ à la \unsure{page|p.} 626 et suivantes).
Dans ce groupe on a défini le sous-groupe, qui était noté
$M_{\rho,\sigma}^{\unsure{\natural|\sharp}}$
\deleted{\unsure{de}}\inserted{\unsure{à}} la p.\ $626'$,
et que je vais plutôt noter maintenant
$M_{\rho_0,\sigma_0}^{\unsure{?|!}}$, défini comme

\margin{686}

% Printed archive footer.
\begin{center}
686\\
\texttt{https://grothendieck.umontpellier.fr}\\
Université de Montpellier, UMR CNRS 5149, Case courrier 051, Place Eugène Bataillon\\
34095 Montpellier cedex 5 - France
\end{center}


\scanpage{688}

\[
\deleted{M_{\unsure{\rho_0|\rho_0,\sigma_0}}^{\unsure{\natural|\sharp}}}
\qquad
\deleted{M_{\rho_0,\sigma_0}^{?}=\unsure{Z_{\rho_0}^{!}}}
\]

\[
(122)\qquad
\deleted{M_{\unsure{\rho_0|\rho_0,\sigma_0}}}
M^?=N_{\unsure{\rho|\rho_0}}^?
\times_{\Gamma_0}Z_{\sigma_0}^!\times_{\Gamma_0}M
\]
où, on rappelle que

\diagram{
\begin{tikzcd}[column sep=small,row sep=small]
&
N_{\rho_0}=\operatorname{Norm}_M(\unsure{L}_{\rho_0}),\quad
N_{\unsure{\sigma|\sigma_0}}=\operatorname{Norm}_M(L_{\sigma_0})
\\
(123)\quad\Biggl\{&
\deleted{\unsure{Z_{\rho_0}}}\qquad
Z_{\unsure{\sigma|\sigma_0}}
=\operatorname{Centr}_M(L_{\sigma_0})
=\operatorname{Centr}_M(\sigma_0)
\\
&
\deleted{N}\inserted{\unsure{Z}}_{\sigma_0}^!
=Z_{\sigma_0}\cap M_0^!
\left(=N_{\sigma_0}\cap M_0^!\right)\subset N_{\sigma_0}
\\
&
N_{\rho_0}^?=Z_{\rho_0}^!\cdot\mu_{\unsure{\tau_1|\tau'}}
\subset N_{\rho_0}\unsure{)|,}\quad
\text{où }Z_{\rho_0}^!=Z_{\rho_0}\cap M^!\ )
\arrow[uu,bend left=45,start anchor=west,end anchor=west]
\end{tikzcd}
}

de sorte qu'on a bien
\[
M^?\subset N_{\rho_0}\times_{\gamma_0}N_{\sigma_0}
\times_{\Gamma_0}M .
\]
\margin{\unsure{indépendant|indépendamment} du choix d'un
$\unsure{\mathcal G}^{\unsure{\natural|\sharp}}$}

Ceci rappelé, on va poser
\begin{samepage}
\[
(124)
\]
\diagram{
\begin{tikzcd}[column sep=small,row sep=small,nodes in empty cells]
M_{\rho_0,\sigma_0}^{\unsure{\natural|\sharp}
 \deleted{\unsure{!|\natural}}?}
\arrow[r,phantom,"\subset" description]\arrow[d,equal]&
M_{\rho_0,\sigma_0}^{\unsure{\natural|\sharp}}
\\
M_{\rho_0,\sigma_0}^{\unsure{\natural|\sharp}}\cap M^?\arrow[d,equal]&
\begin{gathered}
\text{(\unsure{intersection}}\\
\text{dans }N_{\rho_0}\times_{\Gamma_0}N_{\sigma_0}
\times_{\Gamma_0}M\text{)}
\end{gathered}
\\
\deleted{\left(
\unsure{N|M}_{\rho_0}^{\unsure{!|?}}
\times_{\Gamma_0}N_{\sigma_0}^{!}
\times_{\Gamma_0}M^{\unsure{\natural|\sharp}}
\right)}
&
\end{tikzcd}
}
\end{samepage}
et considérons \unsure{l'hom|l'hom.} induit par (119)
\begin{samepage}
\[
(125)
\]
\diagram{
\begin{tikzcd}
M_{\rho_0,\sigma_0}^{\unsure{\natural|\sharp}?}
\arrow[r]\arrow[d,hook]&
M^{\unsure{\natural|\sharp}}\arrow[d,hook]\\
M^?\arrow[r]&M
\end{tikzcd}
}
\end{samepage}
.

\margin{687}

% Printed archive footer.
\begin{center}
687\\
\texttt{https://grothendieck.umontpellier.fr}\\
Université de Montpellier, UMR CNRS 5149, Case courrier 051, Place Eugène Bataillon\\
34095 Montpellier cedex 5 - France
\end{center}


\scanpage{689}
\mspage{647}

On a vu (p.\ $626'$, (55)) que $M^?\longrightarrow M$ fait de $M^?$
une extension de $M$ de noyau $\mu_\rho$.
Donc le noyau de $M_{\rho_0,\sigma_0}^?\longrightarrow
M^{\unsure{\natural|\sharp}}$ est égal à
$\mu_\rho\cap M_{\rho_0,\sigma_0}^?$
or il est clair que l'on a
$\mu_\rho\subset M_{\rho_0,\sigma_0}^{\unsure{\natural|\sharp}}$
\inserted{cf (\unsure{119|114})},
donc $\mu_\rho\subset
M_{\rho_0,\sigma_0}^{\unsure{\natural|\sharp}?}$.

On trouve donc une \deleted{\unsure{suite ex.}}\inserted{hom}
de suites exactes
\begin{samepage}
\[
(126)
\]
\diagram{
\begin{tikzcd}[column sep=small,row sep=small]
&&\inserted{M_{\rho_0,\sigma_0}^{\unsure{\natural|\sharp}}}
\arrow[d,phantom,"\cup" description]&&\\
1\arrow[r]&\mu_\rho\arrow[r]\arrow[d,equal]&
M_{\rho_0,\sigma_0}^{\inserted{\unsure{\natural|\sharp}}?}
\arrow[r]\arrow[d,hook]&
M^{\unsure{\natural|\sharp}}\arrow[d,hook]&\\
1\arrow[r]&\mu_\rho\arrow[r]&M^?\arrow[r]&M\arrow[r]&1
\end{tikzcd}
}
\end{samepage}
et il reste à voir si \unsure{l'}\deleted{\unsure{flèche}}\inserted{hom}
\[
(\infty)\qquad
M_{\rho_0,\sigma_0}^?\twoheadrightarrow M^{\unsure{\natural|\sharp}}
\]
est un épimorphisme, et sinon, \unsure{caractériser} son image.
\deleted{\unsure{Considér|Considérons}}
Notons que \inserted{la définition\unsure{,} (124) donne aussi}
\deleted{\unsure{cet} \unsure{hom} \unsure{s'insère} dans un
\unsure{diagramme}}
compte tenu des expressions (118) pour
\deleted{\unsure{de|des} suites exactes}
$M_{\rho_0,\sigma_0}^{\unsure{\natural|\sharp}}$ et (122) pour $M^?$ :
\[
(127)\qquad
M_{\rho_0,\sigma_0}^{\unsure{\natural|\sharp}?}
\simeq N_{\rho_0}^{\unsure{\natural|\sharp}?}
\times_{\deleted{\unsure{\gamma_{\rho_0,\sigma_0}}}
 \inserted{\Gamma_0^{\unsure{\natural|\sharp}}}}
\deleted{N}\inserted{\unsure{Z}}_{\sigma_0\,0}^{\unsure{\natural|\sharp}!}
\times_{\deleted{\unsure{\gamma_{\rho_0,\sigma_0}}}
 \inserted{\Gamma_0^{\unsure{\natural|\sharp}}}}
M^{\unsure{\natural|\sharp}}
\]
\margin{
\diagram{
\begin{tikzcd}[column sep=small,row sep=small]
\gamma_0\arrow[r,phantom,"\subset" description]\arrow[d,equal]&
\widetilde{\Gamma}_{\mathbb Q}\\
\operatorname{Im}(M^{\unsure{\natural|\sharp}})&
\end{tikzcd}
}
}
où on a posé
\[
(128)\qquad
\left\{
\begin{aligned}
N_{\rho_0}^{\unsure{\natural|\sharp}?}
&=N_{\rho_0}^{?}\cap N_{\rho_0}^{\unsure{\natural|\sharp}}
=N_{\rho_0}^{\unsure{?|!}}\cap M^{\unsure{\natural|\sharp}}\\
\deleted{N}\inserted{\unsure{Z}}_{\sigma_0\,0}^{\unsure{\natural|\sharp}!}
&=Z_{\sigma_0\,0}^!\cap N_{\sigma_0}^{\unsure{\natural|\sharp}}
=\deleted{N}\inserted{\unsure{Z}}
_{\sigma_0\,\underset{\deleted{\unsure{?|\natural}}}{0}}^!
\cap M^{\unsure{\natural|\sharp}}
\end{aligned}
\right.
\]

\margin{688}

% Printed archive footer.
\begin{center}
688\\
\texttt{https://grothendieck.umontpellier.fr}\\
Université de Montpellier, UMR CNRS 5149, Case courrier 051, Place Eugène Bataillon\\
34095 Montpellier cedex 5 - France
\end{center}


\scanpage{690}

et l'hom $(\unsure{\infty|\alpha})$ n'est autre que la troisième projection du produit fibré.
Donc l'image de $(\unsure{\infty|\alpha})$ n'est autre que l'image inverse, dans $M^{\unsure{\natural|\sharp}}$ par $M^{\unsure{\natural|\sharp}}\longrightarrow\gamma_0=\gamma_{\deleted{\unsure{\natural|\sharp}}\rho_0,\sigma_0}$ de l'image de
\[
 (\unsure{\infty\infty|\alpha\alpha})\qquad
 N_{\rho_0}^{\deleted{\unsure{\natural|\sharp}}\inserted{\unsure{\natural|\sharp}?}}
 \mathop{\times}_{\deleted{\unsure{\gamma_{\rho_0,\sigma_0}}}\inserted{\Gamma_0^{\unsure{\natural|\sharp}}}}
 N_{\sigma_0}^{\unsure{\natural|\sharp}}{}_{0}^{!}
 \longrightarrow
 \deleted{\unsure{\gamma_{\rho_0,\sigma_0}}}\inserted{\Gamma_0^{\unsure{\natural|\sharp}}}
 (\subset\Gamma_0)
\]
Il reste donc à montrer que cet hom est surjectif, ou encore que les hom
\diagram{\(\left\{\vcenter{\hbox{\begin{tikzcd}[row sep=large,column sep=huge]
 N_{\rho_0}^{\unsure{\natural|\sharp}?} \arrow[r] & \deleted{\unsure{\gamma_{\rho_0,\sigma_0}}}\inserted{\Gamma_0^{\unsure{\natural|\sharp}}} \\
 N_{\sigma_0}^{\unsure{\natural|\sharp}}{}_{0}^{!} \arrow[r] & \Gamma_0^{\unsure{\natural|\sharp}}
\end{tikzcd}}}\right.\)}
sont surjectifs (le premier sera alors une ext. de $\Gamma_0^{\unsure{\natural|\sharp}}$ par $\mu_\rho$, le deuxième un iso). Soit donc $u\in\Gamma_0^{\unsure{\natural|\sharp}}$, \deleted{défini par un}\inserted{provenant de} $u_{\unsure{0}}\unsure{\in}M^{\unsure{\natural|\sharp}}\deleted{\unsure{_{\rho,\sigma}}}\simeq G_{\rho,\sigma}^{\unsure{\natural|\sharp}}$ satisfaisant donc aux conditions a) b) c) d) de (107). On voit donc qu'il existe $\beta_0\in\unsure{\mathcal G}_0^{\unsure{\natural|\sharp}}\inserted{\subset M^{\unsure{\natural|\sharp}}}$ tel que $\beta^{-1}u_0\in N(\sigma_0)$.

\margin{Ici j'ai évité les vérifications, pensant y revenir dans les pages suivantes, il s'avère que ce n'est pas urgent -- il importait plutôt de développer les cas particuliers ``\unsure{tour|tours} de Fermat'' pour préciser les idées \inserted{des formulations \unsure{générales} des $M_{\rho,\sigma}$\ldots.}}

\margin{689}

% Printed archive footer.
\begin{center}
689\\
\texttt{https://grothendieck.umontpellier.fr}\\
Université de Montpellier, UMR CNRS 5149, Case courrier 051, Place Eugène Bataillon\\
34095 Montpellier cedex 5 - France
\end{center}


\scanpage{691}
\mspage{648}

$9^{\circ})$ \uline{Digression de formalisme des groupes}.

\textcircled{a}

Soient
\[
 (1)\qquad
 \text{\uline{$G$}},\quad
 \text{\uline{$\unsure{\mathcal G}\subset G$}},\quad
 \text{\uline{$\unsure{N}\overset{\varphi_N}{\hookrightarrow}G$}},\quad
 \inserted{\text{d'où}}\quad
 SN=\deleted{\unsure{N\cap\mathcal G}}\inserted{\varphi^{-1}(\unsure{\mathcal G})}\subset N
\]
la situation d'un groupe $G$, sous-groupe invariant $\unsure{\mathcal G}$, et sous-groupe $N$ tel que
\[
 (\unsure{\alpha})\qquad
 N\xrightarrow{\text{épi}}G/\unsure{\mathcal G}\deleted{\unsure{\overset{\mathrm{dfn}}{=}\gamma}},
 \qquad\text{i.e.}\quad N/SN\xrightarrow{\sim}G/\unsure{\mathcal G}
\]
\deleted{\unsure{soit épi}} On va reconstituer cette situation à partir de la suivante
\[
 (2)\qquad
 \unsure{\mathcal G},\quad N,\quad
 SN\subset N\deleted{\unsure{\operatorname{Aut}}},\quad
 N\xrightarrow{\deleted{\unsure{\varphi}}\inserted{\Theta}}\operatorname{Aut}(\unsure{\mathcal G}),
 \quad\inserted{SN\xrightarrow[\varphi=\varphi_{SN}]{\deleted{\unsure{\varphi}}}\unsure{\mathcal G}}
\]
où $\unsure{\mathcal G},N$ sont des groupes, $SN$ un sous-groupe distingué de $N$, $N\longrightarrow\operatorname{Aut}(\unsure{\mathcal G})$ un hom de groupes, i.e. une \uuline{action} de $N$ sur $\unsure{\mathcal G}$; enfin $SN\longrightarrow\unsure{\mathcal G}$ un hom de groupes. On suppose que \deleted{l'action}
\deleted{que $SN\longrightarrow\unsure{\mathcal G}$ commute}
\[
 (3)\quad
 \left\{\begin{aligned}
 &\inserted{\text{a) l'action}}\text{ de }SN\text{ sur }\unsure{\mathcal G}\text{ \unsure{induit|induite} par celle de }\unsure{\mathcal G|G},\\
 &\text{soit aussi celle déduite de }SN\longrightarrow\unsure{\mathcal G},\\
 &\text{via l'op. adjointe de }\unsure{\mathcal G},\\
 &\inserted{\text{b) }SN\longrightarrow\unsure{\mathcal G}\text{ est compatible aux actions de }N.}
 \end{aligned}\right.
\]
et \deleted{\unsure{Il est|Je vois}} \deleted{\unsure{Il est|Je vois}} clair comment (1) définit (2), montrons comment un syst. (2) permet de reconstituer un syst. (1).

\margin{690}

% Printed archive footer.
\begin{center}
690\\
\texttt{https://grothendieck.umontpellier.fr}\\
Université de Montpellier, UMR CNRS 5149, Case courrier 051, Place Eugène Bataillon\\
34095 Montpellier cedex 5 - France
\end{center}


\scanpage{692}

On part donc de \unsure{(2)}, on pose
\[
 (4)\qquad S_0G=N.\unsure{\mathcal G}\qquad(\text{\unsure{Ch direct|pr. direct}}),
\]
on définit un hom
\diagram{\((5)\quad\left\{\vcenter{\hbox{\begin{tikzcd}[row sep=large,column sep=large]
 SN \arrow[r,hook,"i"] & S_0G \\
 \unsure{u} \arrow[r,mapsto] & \unsure{u}.\varphi(\unsure{u})^{-1}=\varphi(u)^{-1}.u
\end{tikzcd}}}\right.\)}

\margin{\uline{NB} on a \inserted{(par hyp. (a))} pour $u\in SN$, $x\in\unsure{\mathcal G}$
\[
 (*)\quad\left\{u(x)=\underset{\substack{\text{dans }S_0G\\=N.\unsure{\mathcal G}}}{\underbrace{uxu^{-1}}}=\varphi(u)x\varphi(u)^{-1}\right.
\]
donc si $x=\varphi(u)$, on a
\[
 u(x)=uxu^{-1}=x
\]
i.e. $x=\varphi(u)$ et $u$ commutent.}

Vérifions que c'est un hom
\deleted{\unsure{on}}
\[
 i(\unsure{u'}u)=(u'u)\varphi(u'u)^{-1}
\]
\[
\begin{aligned}
 i(u')i(u)&=u'\varphi(u')^{-1}u\varphi(u)^{-1}\\
 &\deleted{=u'\varphi(u')^{-1}u\varphi(u')}\\
 &=u'u\underset{\substack{=\varphi(u^{-1})\varphi(u'^{-1})\varphi(u)\\\text{par a)}}}{\underbrace{(u^{-1}\varphi(u'^{\unsure{-1}})u)}}\varphi(u^{-1})^{\deleted{\unsure{-1}}}\\
 &=u'u\varphi(u^{-1}u'^{-1}\deleted{uu^{-1}})=(u'u)\varphi((u'u)^{-1})=(u'u)\varphi(u'u)^{-1}\qquad\text{ok.}
\end{aligned}
\]

On n'a pas encore utilisé le fait que soit distingué dans $N$, \inserted{et l'hyp. (b)} ceci implique que l'image de \unsure{(4)|(5)} est distinguée dans $S_0G$, car \deleted{elle est} on a, si $a\in N$, $u\in SN$
\[
 \deleted{\unsure{(6)(6)}}\qquad\deleted{\unsure{(*)}}\quad
 a(u.\varphi(u)^{-1})a^{-1}
 =\underbrace{aua^{-1}}_{\in SN}.\underset{=\varphi(au^{-1}a^{-1})=\varphi(au^{-1}a^{-1})^{-1}}{\underbrace{a\varphi(u^{-1})a^{-1}}}
\]
i.e. l'hom (5) \deleted{\unsure{est}} commute aux actions de $N$ sur $SN$ et sur $S_0G\supset N$ -- donc l'image de $SN$ est normalisée par $N\subset S_0G$. Il faut encore vérifier qu'elle est normalisée également par $\unsure{\mathcal G}$, \deleted{soit donc $x\in\unsure{\mathcal G}$, on a}\inserted{mais on a mieux :}
\[
 \deleted{xu\varphi(u^{-1})x^{-1}=(xux^{-1})(x\varphi(u^{-1})x^{-1}}
\]
L'image de $SN\xrightarrow{i}S_0G$ \uline{commute} à $\unsure{\mathcal G}$, i.e.

\margin{691}

% Printed archive footer.
\begin{center}
691\\
\texttt{https://grothendieck.umontpellier.fr}\\
Université de Montpellier, UMR CNRS 5149, Case courrier 051, Place Eugène Bataillon\\
34095 Montpellier cedex 5 - France
\end{center}


\scanpage{693}
\mspage{649}

$\forall u\in SN$, $u\varphi(u)^{-1}$ induit l'opération triviale sur $\unsure{\mathcal G}$ --- ce qui est en effet trivial par (3) a).

On peut donc poser
\[
(6)\qquad  G\overset{\unsure{\mathrm{d\acute ef}|\mathrm{dfn}}}{=}S_0G/i(SN),
\]
on trouve alors un diagramme commutatif

\begin{samepage}
\noindent (7)\quad
\diagram{\begin{tikzcd}[column sep=small,row sep=small]
& N \arrow[dr,"\varphi_N"] & \\
SN \arrow[ur,hook] \arrow[dr,"\varphi_{SN}"'] & & G \\
& \unsure{\mathcal G} \arrow[ur,hook] &
\end{tikzcd}}
\end{samepage},

on voit que l'image de $\unsure{\mathcal G}$ dans $G$ est un ss-groupe invariant, et que l'\unsure{on} a un isom. sur les quotients:
\[
(8)\qquad  N/SN\xrightarrow{\sim}G/\unsure{\mathcal G},
\]
ce qui implique que l'on a aussi un iso
\[
(9)\qquad  \operatorname{Ker}(\varphi_{SN}:SN\longrightarrow\unsure{\mathcal G})
\xrightarrow{\sim}\operatorname{Ker}(N\xrightarrow{\varphi_N}G);
\]
\deleted{donc $\varphi_N$}

\begin{samepage}
\noindent (10)\quad
\diagram{\(\begin{array}{l}
\varphi_N:N\longrightarrow G\text{ est un plongt}\\
\unsure{\Updownarrow}\\
\varphi_{SN}:SN\longrightarrow\unsure{\mathcal G}\text{ est un plongement.}
\end{array}\)}
\end{samepage}

On trouve ainsi des \uline{équiv. de catégories}, inverses l'une de l'autre, entre \inserted{\unsure{les cat.}} des

\margin{692}

% Printed archive footer.
\begin{center}
692\\
\texttt{https://grothendieck.umontpellier.fr}\\
Université de Montpellier, UMR CNRS 5149, Case courrier 051, Place Eugène Bataillon\\
34095 Montpellier cedex 5 - France
\end{center}


\scanpage{694}

situations (1), et celle des situations (2).

\textcircled{b}
Considérons maintenant une situation
\[
(11)\qquad 
\left\{\begin{array}{l}
G,\quad \unsure{\mathcal G}\subset G,\quad (\varphi_i:N_i\longrightarrow G)_{i\in I}\\[4pt]
\text{d'où des }SN_i=\varphi_i^{-1}(\unsure{\mathcal G})
\end{array}\right.
\]
sous la condition
\[
(12)\qquad 
\forall i\in I,\quad \varphi_i:N_i\longrightarrow G/\unsure{\mathcal G}=\gamma\quad\text{épi}
\]
\unsure{peut} on la reconstituer à partir de
\[
(13)\qquad 
\unsure{\mathcal G},\ (N_i,SN_i,\unsure{\varphi_i},\Theta_i)_{i\in I}
\]
où \inserted{$\forall i\in I,$} $(N_i,SN_i,\varphi_i,\Theta_i)$ satisfait les conditions (3) a) et b) --- \unsure{où|et} bien sûr $\forall i\in I$
\[
(14)\qquad 
\deleted{\unsure{\varphi_i}}\inserted{\Theta_i}:N_i\longrightarrow\operatorname{Aut}(\unsure{\mathcal G}),\qquad
\varphi_i:SN_i\longrightarrow\unsure{\mathcal G}.
\]

Il est tentant de reconstituer $G$ à l'aide \unsure{de|des} $\unsure{\mathcal G}$ et des $N_i$, comme une sorte de ``somme amalgamée'' relative à:

\noindent \deleted{\unsure{(15)}}\quad
\diagram{\begin{tikzcd}[column sep=small,row sep=small]
& \unsure{\mathcal G} \arrow[dr,dashed] & \\
SN_i \arrow[ur,"\varphi_i"] \arrow[dr,hook,"\mathrm{plongt}"'] & & G \\
& N_i \arrow[ur,dashed] &
\end{tikzcd}}

\margin{693}

% Printed archive footer.
\begin{center}
693\\
\texttt{https://grothendieck.umontpellier.fr}\\
Université de Montpellier, UMR CNRS 5149, Case courrier 051, Place Eugène Bataillon\\
34095 Montpellier cedex 5 - France
\end{center}


\scanpage{695}
\mspage{650}

\deleted{\uline{Proposition} \unsure{Partant|Partir} de (11), d'où (13)}

% small x-shaped marks follow "dire," and $\Theta_i$ in this sentence (reference marks?); not transcribed as text.
Mais il est plus vraisemblable de dire, que pour tt $i$, on reconstitue \unsure{via|ici} $N_i,SN_i,\varphi_i,\Theta_i$ la situation (pour $i$ fixé)

\noindent \deleted{\unsure{(17)}}\inserted{\unsure{(15)}}\quad
\diagram{\begin{tikzcd}
\deleted{\unsure{\varphi_i}}N_i \arrow[r] & G_i(\unsure{\simeq}G) \\
SN_i \arrow[u,hook] \arrow[r] & \unsure{\mathcal G} \arrow[u,hook]
\end{tikzcd}}

et qu'on reconstitue (11) lui-$\overline{\mathrm m}$ en se donnant un \uline{syst. transitif} \uline{d'isom} \uline{entre les} $G_i$ $(i\in I)$, \uline{induisant} \uline{l'identité sur} $\unsure{\mathcal G}$. Une telle donnée implique en premier lieu la donnée d'un syst. transitif d'iso entre les
\[
(16)\qquad 
\gamma_i\overset{\unsure{\mathrm{d\acute ef}|\mathrm{dfn}}}{=}N_i/SN_i
\]
--- qu'on \unsure{identifiera} donc à un $\overline{\mathrm m}$ groupe
\[
(17)\qquad 
\begin{gathered}
\gamma,\text{ groupe quotient commun}\\
\text{des }N_i\text{ (par les }SN_i\text{).}
\end{gathered}
\]
Identifions maintenant les $G/\unsure{\mathcal G}_i=\gamma_i$ à $\gamma$
\[
(18)\qquad G_i/\unsure{\mathcal G}\simeq\gamma,
\]
et formons le produit fibré

\margin{694}

% Printed archive footer.
\begin{center}
694\\
\texttt{https://grothendieck.umontpellier.fr}\\
Université de Montpellier, UMR CNRS 5149, Case courrier 051, Place Eugène Bataillon\\
34095 Montpellier cedex 5 - France
\end{center}


\scanpage{696}

\[
(19)\qquad 
\unsure{g}=\underset{\text{(produit fibré sur }\gamma\text{)}}{\prod_{i\in I}\!{}_{\gamma}G_i}\ \supset\ \unsure{\mathcal G}^I
\]
\margin{$\unsure{g}$ est extension de $\gamma$ \unsure{par} $\unsure{\mathcal G}^I$}

La donnée d'un syst. transitif d'iso entre les $G_i$, induisant l'identité sur les sous-groupes $\unsure{\mathcal G}$, et compatible avec le système \deleted{\unsure{d'iso|du}} \unsure{transitif} d'iso. déjà envisagé entre les $\gamma_i=N_i/SN_i\simeq G_i/\unsure{\mathcal G}$, équivaut alors à la donnée d'un sous-groupe $G$ de $\unsure{g}$, satisfaisant les conditions
\[
(20)\qquad 
\left\{\begin{array}{ll}
\text{a)}&G\cap\unsure{\mathcal G}^I=\text{diagonale }\delta(\unsure{\mathcal G})\text{ de }\unsure{\mathcal G}^I\\[5pt]
\text{b)}&\text{La projection }G\longrightarrow\gamma\text{ est épi}
\end{array}\right.
\]

En d'autres termes, on a construit, à l'aide de (13) et du syst. transitif d'iso. entre les $\gamma_i$, une extension $\unsure{g}$ de leur ``valeur commune'' $\gamma$, par $\unsure{\mathcal G}^I$ --- et la structure supplémentaire, correspondant à la donnée \unsure{(11)}, est la ``restriction'' de cette extension en une extension de $\gamma$ par le sous-groupe diagonal de $\unsure{\mathcal G}^I$.

\margin{695}

% Printed archive footer.
\begin{center}
695\\
\texttt{https://grothendieck.umontpellier.fr}\\
Université de Montpellier, UMR CNRS 5149, Case courrier 051, Place Eugène Bataillon\\
34095 Montpellier cedex 5 - France
\end{center}
